% Starter template for an accessible (tagged, PDF/UA-1) set of notes.
% Compile with LuaLaTeX on TeX Live 2025 or newer (Overleaf: Menu > Compiler = LuaLaTeX,
% TeX Live version = 2025 or 2026). Run it twice so cross-references resolve.
\DocumentMetadata{
  lang        = en-US,
  pdfstandard = ua-1,
  pdfversion  = 1.7,
  tagging     = on,
}
\documentclass[11pt]{article}

\usepackage{amsmath}
\usepackage{unicode-math}
\setmathfont{Latin Modern Math}
\usepackage{graphicx}
\usepackage[margin=1in]{geometry}
\usepackage{hyperref}

\hypersetup{
  pdftitle  = {Newton's Laws and Friction},   % shown as the document title by screen readers
  pdfauthor = {Your Name},
  hidelinks,
}

\begin{document}

\begin{center}
  {\LARGE\bfseries Newton's Laws and Friction}\\[6pt]
  {\large Physics 101 \quad Your Name}
\end{center}

% Use \section and \subsection for every heading. Never fake a heading with \textbf.
\section{Newton's second law}

The net force on an object equals its mass times its acceleration:
\begin{equation}
  \vec F_{\text{net}} = m\,\vec a .
  \label{eq:newton2}
\end{equation}
Equation~\eqref{eq:newton2} holds in an inertial reference frame.

\section{Friction}

\subsection{Static and kinetic friction}

Use real list environments for lists:
\begin{itemize}
  \item Static friction satisfies $f_s \le \mu_s N$.
  \item Kinetic friction satisfies $f_k = \mu_k N$.
\end{itemize}

% Every image needs alt text that says what the picture shows physically.
\begin{center}
  \includegraphics[width=0.4\linewidth,
    alt={A box resting on a horizontal floor. An arrow labeled F pushes the box
         to the right; an arrow labeled f, for friction, points to the left;
         an arrow labeled N points up and an arrow labeled mg points down.}]
    {example-image}
\end{center}

% Tables: declare the header row before the tabular so it is tagged as a header.
\tagpdfsetup{table/header-rows={1}}
\begin{center}
\begin{tabular}{lcc}
  Surfaces & $\mu_s$ & $\mu_k$ \\
  Rubber on dry concrete & 1.0 & 0.8 \\
  Steel on steel & 0.74 & 0.57 \\
\end{tabular}
\end{center}

\end{document}
